\section{Introduction}

An analytics team that adopts a semantic layer is making a testable bet: that
encoding business definitions once, in governed objects a query engine compiles
rather than a language model improvises, makes a downstream agent answer new
questions more accurately than the same agent writing SQL against a raw schema.
That bet is hard to evaluate from existing comparisons, which typically move
several things at once and leave the harness, the question selection, and the
number of scored runs unreported.

Three separable things get bundled into such a comparison. \emph{Access}: whether
the business knowledge needed to answer the question is available at all. A
question about ``active high-value accounts'' cannot be answered from column
names, and no architecture repairs a missing definition. \emph{Representation}:
whether that knowledge is a paragraph of prose the model reads or a structured
object with declared dependencies, grain, and join paths. \emph{Enforcement}:
whether the system may ignore the representation and emit whatever SQL it likes,
or must compile its intent through the governed layer and pass its validation.
The three carry very different costs, so a comparison that moves all three at
once cannot attribute the difference to any of them.

This paper preregisters an experiment that holds them apart. We evaluate the Omni
semantic layer on LiveSQLBench Large-v1, a BIRD-family text-to-SQL benchmark
whose 18 PostgreSQL databases ship the material a semantic layer is supposed to
encode: a hierarchical knowledge base (HKB) of 1,090 business definitions
connected by 945 declared dependency edges, 344 of which point at another derived
entry, with multi-hop resolution required in every database. A benchmark in which
business knowledge composes is one where the difference between a prose blob and
a dependency-aware model can actually show up.

The design is a four-rung ladder, frozen together and evaluated in one sealed
event. C1 gives a competent agent the raw schema and lets it write SQL. C2 adds
programmatic search over the raw HKB, changing access alone. C3 replaces the raw
HKB with a searchable export of the compiled semantic model, changing
representation while the agent still writes its own SQL. C4 is the
production-governed system, where semantic compilation and validation are
enforced. The primary product endpoint is C4 mean one-shot execution accuracy
across three independent repetitions of the 101 sealed questions; the primary
comparative endpoint is the paired C4 minus C1 difference on the same trials.
Rung-to-rung contrasts are prespecified but exploratory, and we state in advance
which can carry a causal reading.

Several commitments exist because of failure modes documented in the evaluation
literature rather than anything specific to this product. All 1,212 generations
complete before any output is scored, so no held-out outcome can reach a prompt,
a retrieval configuration, or a retry policy. The development partition is split
again into a development partition that the protocol permits reusing and a
metered partition available at most ten times behind a signing boundary. The
executed final candidate uses neither question-level supervision nor a metered
checkpoint. Two scorers are
frozen before any gold is released, one reproducing the official evaluator's
known lossy behavior and one repairing it, with neither selectable after results
are seen. Every attempt carries a telemetry envelope in which an unobserved token
count is null with a named reason, never zero. Every condition's harness is
disclosed, because controlled studies place the scaffold's effect on agent
accuracy anywhere from negligible to decisive.

The paper reports no accuracy, and that is the current state of the work: the
split and protocol are frozen, custody tooling and both scorers are implemented,
and the capture verification gate that must precede any scaled run has not been
passed. Publishing the design before the numbers exist is the point. The
contributions are a public-artifact reconnaissance of LiveSQLBench Large-v1
including its knowledge-graph structure and the exact behavior of its pinned
evaluator; a four-condition ablation separating access, representation, and
enforcement, with an explicit account of which contrasts survive the absence of
model parity; a custody and freeze protocol under which supervised development on
231 questions cannot contaminate 101 sealed ones; a machine-checkable development
control plane in which hypotheses precede changes and rejected experiments cannot
be rewritten out of the record; and a telemetry contract that keeps correctness
out of the immutable generation record and missing measurements distinguishable
from measured zeros.
